\subsection{Runtime Performance}\label{subsec:runtime_perf}

For the evaluation of the computational efficiency of the implemented approaches, we examine the runtime impact of the optimizations described in Section~\ref{subsec:optimizations}.

\paragraph{Optimization Impact.}
Table~\ref{tab:table_bataineh} summarizes the progressive runtime improvements obtained for the method of Bataineh et al.~\cite{Bataineh2011}.
Starting from a naive baseline implementation with explicit sliding-window computations, successive optimizations significantly reduced the execution time.

% Table 1: Compares different optimizations for Bataineh's approach
\begin{table}[t]
    \begin{tabular}{lcccc}
                  & Base & Single-Pass & Integral Image & OpenMP \\
        \hline
        Time (ms) & 1330 & 748         & 188            & 195
    \end{tabular}
    \caption{Runtime comparison of optimization stages for Bataineh's method on the medium-resolution image.}
    \label{tab:table_bataineh}
\end{table}

The results demonstrate that algorithmic improvements, particularly the use of integral images, provide the largest performance gain.
Overall, the optimized implementation achieves an approximate $7\times$ speedup compared to the baseline.
Interestingly, the OpenMP parallelized version does not yield further improvement, suggesting that the computation becomes limited by memory access overhead and parallelization costs rather than arithmetic operations.

\paragraph{High-Resolution Image Runtime.}
To better distinguish the performance characteristics of the different approaches, we also benchmarked all optimized methods on a high-resolution image of size $6700 \times 4467$ pixels.
The resulting runtimes are reported in Table~\ref{tab:table_compare}.

% Table 2: Comparing the different approaches on a 4K picture
\begin{table}[t]
    \begin{tabular}{lc}
        Approach & Time (s) \\
        \hline
        Niblack  & 6.30 \\
        Sauvola  & 6.44 \\
        Nick     & 6.34 \\
        Bataineh & 6.39 \\
        Bradley  & 4.55 
    \end{tabular}
    \caption{Running times of the optimized approaches on the high-resolution image.}
    \label{tab:table_compare}
\end{table}

All classical window-based methods exhibit similar runtime behavior, as they require both mean and standard deviation statistics for each pixel.
Bradley's approach achieves the lowest runtime, which can be attributed to its reliance on mean-based thresholding only, combined with the efficiency of integral image evaluation.
